\section{Limitations}

The primary external cohort contains 75 participants. Differences between HBN
and MIPDB in recruitment, age distribution, hardware, and recording procedures
are combined in the observed transfer performance; this design cannot identify
a single source of error. The minimum sample-size threshold did not guarantee
precise estimates. Intervals that include zero leave room for both useful gains
and worse performance, and no equivalence margin was specified.

The main comparison covers one REVE checkpoint, four heads, and ten seeds under
common training settings. Equal budgets do not guarantee that each head was
optimally tuned. The training-size analysis covers three nested subsets, and
the layer-wise analysis covers four depths with separately fitted linear heads.
Both reuse MIPDB after the primary evaluation. These results do not establish a
scaling law or a general ranking of encoders, layers, or heads.

The ds006780 extension uses a separate cohort but was motivated by earlier
results. Its reference was undocumented and its electrode positions were
approximated by a canonical BioSemi montage. The source dataset includes a
sibling group. Family links were not retained in the aggregate evidence, so we
could not establish whether related participants entered the selected sample.
The bootstrap treats participants as independent; any unaccounted family
clustering could make its intervals too narrow. A family-aware sensitivity
analysis would require the original cohort metadata.

The HBN source manifest records age and split membership but not sex or
diagnostic group, limiting the assessment of sampling imbalance. We recovered
the selected ds006780 sample description from public metadata and verified its
participant and age hashes. These checks do not recover the original prediction
vectors or execution caches. Independent recalculation of participant-level
metrics requires those artifacts or a rerun using the original data.

MIPDB participants outside the HBN training-age range were excluded from the
primary analysis. REVE's published pretraining inventory includes HBN releases,
so HBN is not an encoder-unseen evaluation. Although the inventory does not name
MIPDB, we could not independently verify complete absence of pretraining overlap.
We make no encoder-unseen claim for ds006780 either. Finally, predicting
chronological age does not validate a clinical brain-age biomarker.
